% !PRINTED 71
\weberpage{71}

In the representation \S~20, (1) of \(\Delta\), the rows are characterized by the upper indices, the columns by the lower indices. Hence this theorem may also be expressed as follows:

\gesperrt{II. A determinant is not changed when the rows are made columns and the columns are made rows.}

If, in all arrangements \(\mathfrak A,\mathfrak A',\mathfrak A'',\ldots\) of the \(n\) elements, we interchange any two elements with one another, the totality of these arrangements remains unchanged, but every arrangement of the first kind passes into an arrangement of the second kind, and conversely. If, therefore, in the terms \(M,M',M'',\ldots\) of which \(\Delta\) is composed, any two lower indices are interchanged, every term with positive sign passes into another which in \(\Delta\) bore the negative sign, and conversely; therefore \(\Delta\) changes its sign. With the help of II follows the theorem:

\gesperrt{III. If in \(\Delta\) two lower or two upper indices are interchanged with one another, the determinant changes only its sign.}

Expressed somewhat differently: \gesperrt{if two rows or two columns are interchanged with one another, the determinant changes only its sign}; and from this more generally:

\gesperrt{IV. If in a determinant the rows or the columns are permuted, the absolute value is not changed, and the sign is not changed or passes into the opposite sign according as the applied permutation belongs to the first or second kind.}

From III one obtains the following fundamental theorem:

\gesperrt{V. If in two rows or in two columns the terms standing in the same place are equal to one another (more briefly, if two rows are equal to one another), then the determinant has the value zero.}

For the interchange of the two rows changes the sign by III, but on the other hand, since both rows are identical,
%CONT>>
